\documentclass{article}
\usepackage[utf8]{inputenc}
\usepackage{geometry}
\usepackage{enumitem}
\usepackage{hyperref}
\usepackage{xcolor}
\usepackage{graphicx}
\usepackage{array}
\usepackage{multirow}
\usepackage{amsmath, amssymb}
\usepackage{chemfig}
\usepackage{mhchem}
\usepackage{tikz}
\usepackage{setspace}
\usepackage{soul}
\usepackage{mathtools}
\usepackage{booktabs}
\usepackage{chemformula}
\usepackage{siunitx}
\usetikzlibrary{positioning, calc}
\usepackage{longtable}
\geometry{margin=1in}
\setlength{\parskip}{0.5em}
\usepackage{cancel}

\begin{document}

\begin{center}
    {\LARGE\textbf{SI Session Planning Form}}
\end{center}

\
\noindent
\begin{flushleft}
\textbf{SI Leader:} Brandon Cobb\\
\textbf{Session Week:} 17\\
\textbf{Session\#:} 17\\
\textbf{Instructor:} Professor Tramontozzi\\
\textbf{Old Plans:}\href{https://macombedu.sharepoint.com/:b:/s/ASC-SupplementalInstructionEmbeddedTutoring/EeSYqCwT_KhMolgJvXbVEjgBOyF_DevvS8BgCm1zPkQQJg?e=QJTt4o}{SharePoint}\\
\textbf{Session Goal:} Students will review and apply core concepts of concentration and solution chemistry, including calculating molarity, mass percent, volume percent, and mass/volume percent, as well as performing dilution calculations and using solubility rules to determine saturation levels and predict precipitation. They will analyze factors that influence solubility, such as temperature and solute–solvent polarity, and explore colligative properties including freezing point depression and boiling point elevation for electrolytes and nonelectrolytes. Students will strengthen their proficiency in stoichiometry by balancing equations, applying atomic masses, chemical formulas, Avogadro’s number, and mole ratios, identifying limiting reactants, and determining theoretical yields. They will classify reaction types—synthesis, decomposition, single and double replacement, combustion, neutralization, and redox—and apply gas laws and kinetic molecular theory to relationships among pressure, volume, temperature, moles, and molecular motion. Students will refine quantitative skills through unit conversions, algebraic manipulation of exact and inexact numbers, correct use of significant figures in multistep calculations, measurement techniques using instruments with fixed uncertainties, proper recording of units, and scientific notation. Additionally, they will deepen their understanding of atomic and subatomic structure, including electron orbitals. By the end of this unit, students will distinguish dissociation from ionization, calculate pH, pOH, ([H_3O^+]), and ([OH^-]), determine percent ionization of weak acids, interpret (K_a) and pKa values, identify conjugate acid–base pairs, analyze polyprotic acids, and predict pH changes caused by salt hydrolysis. They will apply these concepts to assess whether solutions are acidic, basic, or neutral and verify that their calculations are consistent, properly unitized, and presented with correct significant figures.
\end{flushleft}

\vspace{1em}
\section*{\underline{Opener}}
{\centering\large\textbf{\hl{Take Attendance}}\par}\vspace{-0.25\baselineskip}
\noindent
\begin{tabular}{|p{4.2cm}|p{9.8cm}|}
\hline
\textbf{Collaborative Learning Techniques} & Turn to a Partner\\
\hline
\textbf{Learning Strategy} & Ice Breaker (Two truths and a Lie) \\
\hline
\textbf{Time} & 11:30---11:35\\
\hline
\textbf{Materials} & Notes, homework, whiteboard, markers, ready to speak\\
\hline
\textbf{Lecture/Textbook/Old Plans/Slide References} & Class roster \\
\hline
\end{tabular}

\vspace{0.5cm}
{\large\textbf{Definitions}}
\\
\begin{tabular}{|p{0.9\linewidth}|}
\hline
\begin{minipage}[t]{0.9\linewidth}
\begin{enumerate}[label=\alph*.]
    \item\textbf{Turn to a Partner:}  Group members work with a partner on an assignment or discussion topic. This technique works best with group participants who have already been provided with enough background on a subject that they can immediately move to a discussion with their partner without previewing or reviewing concepts. 
   \item\textbf{Ice Breaker:} Each student states two true facts and one false statement related to chemistry concepts (or general academic experiences), and the group guesses which statement is false. This activity encourages peer interaction, sparks discussion, and lightly reviews content in a fun, low-pressure way. 
\end{enumerate}
\end{minipage} \\
\hline
\end{tabular}
\newpage

\noindent
{\large\textbf{Checking for Understanding (questions))}}\\

\begin{tabular}{|p{0.9\linewidth}|}
\hline
\begin{itemize}
    \item What is the primary goal of this activity?

    
    \item How do students know if their solution at a station is correct?
.
    
    \item What should students do if they are unsure about a question?

    
    \item How does moving between stations benefit student learning?

    
    \item What is the importance of showing work on the team slip?

    
    \item How are bonus points awarded during the activity?
n.
    
    \item What is the role of the facilitator at each station?

    
    \item How should teams handle disagreements while solving a question?

    
    \item Why is it important that only verified answers are recorded on the team slip?

    
    \item What is the wrap-up discussion intended to accomplish?
\end{itemize} \\
\hline
\end{tabular}




\section*{\underline{Activity 1}}
\noindent
\begin{tabular}{|p{4.2cm}|p{9.8cm}|}
\hline
\textbf{Collaborative Learning Techniques} & Group Discussion\\
\hline
\textbf{Learning Strategy} & Lecture Review\\
\hline
\textbf{Time} & 11:40---12:05\\
\hline
\textbf{Materials} & Notes, homework, whiteboard, markers, ready to speak, periodic table \\
\hline
\textbf{Lecture/Textbook/Old Plans/Slide References} & Chapters 1-10\\
\hline
\end{tabular}

\vspace{0.5cm}
{\large\textbf{Definitions}}
\\
\begin{tabular}{|p{0.9\linewidth}|}
\hline
\begin{minipage}[t]{0.9\linewidth}
\begin{enumerate}[label=\alph*.]
    \item\textbf{Group Discussion:} One person in the group is asked to present on a topic or review material for the group and then lead the discussion for the group. This person should not be the regular group leader.
   \item\textbf{Lecture Review:} As a group summarize the lecture from the previous class. You may have to provide prompts for the students. For example, "The first concept discussed was Civil Liberties and Public Policy, what did the profressor highlight regarding this?" You may want to ask them to try summarizing without looking at their notes; however, if they are having a difficult time remembering, tell them to refer to their notes.
\end{enumerate}
\end{minipage} \\
\hline
\end{tabular}
\newpage
\noindent
\begin{tabular}{|p{0.9\linewidth}|}
\hline
{\large\textbf{Activity Overview}} \\
\begin{enumerate}[label=\alph*.]
    \item \textbf{Purpose:} Provide students with a collaborative, lecture-integrated review session that strengthens comprehension, encourages active participation, and reinforces key concepts from the most recent lesson. Students work together to analyze prompts, discuss big ideas, and apply them to example scenarios.
    
    \item \textbf{Structure:} 
    \begin{itemize}
        \item Students work in small groups (3–5).
        \item The review proceeds in several timed “rounds,” each focused on a major idea from the lecture (e.g., mechanisms, definitions, graphical reasoning, model interpretation).
        \item Each round involves a shared prompt projected for the class (or distributed as slips), followed by structured group discussion and a brief written response.
        \item After each round, the instructor leads a whole-class debrief before transitioning to the next concept.
    \end{itemize}
    
    \item \textbf{Materials:} 
    \begin{itemize}
        \item Prepared question prompts for each review round.
        \item Group response sheets or whiteboards.
        \item Timer or projected countdown.
        \item Optional visual supports (figures, graphs, chemical structures).
    \end{itemize}
\end{enumerate} \\
\hline
\end{tabular}
\newpage

\noindent
{\large \textbf{Instructions (detailed activity breakdown below) continued...}} \\
\vspace{0.5cm}
\noindent
\begin{tabular}{|p{0.9\linewidth}|}
\hline
{\large\textbf{Facilitation Instructions}} \\
\begin{enumerate}[label=\alph*.]
    \item \textbf{Group Management:} 
    \begin{itemize}
        \item Assign students to balanced groups based on seating or prior performance.
        \item Encourage each group to designate rotating roles (scribe, presenter, checker, timekeeper).
    \end{itemize}

    \item \textbf{Round Interaction:} 
    \begin{itemize}
        \item Present the prompt for the round and start the timed interval.
        \item Groups discuss the question, synthesize reasoning, and record a short, unified response.
        \item Encourage equal participation within groups; visit groups and ask probing questions.
        \item When time is up, instruct groups to finalize responses before the debrief.
    \end{itemize}

    \item \textbf{Instructor Role:}
    \begin{itemize}
        \item Circulate to observe discussions and identify misconceptions.
        \item Use targeted guiding questions to stimulate reasoning.
        \item Collect insights during group work to address during whole-class debriefs.
        \item Maintain pacing across rounds to keep the review session energetic and focused.
    \end{itemize}
\end{enumerate} \\
\hline
\end{tabular}

\vspace{0.5cm}
\noindent
\begin{tabular}{|p{0.9\linewidth}|}
\hline
{\large\textbf{Rules and Expectations}} \\
\begin{enumerate}[label=\alph*.]
    \item All group members must contribute meaningfully to each round.
    \item Responses must reflect consensus, not individual work.
    \item Respectful discourse is expected; challenges to ideas should be evidence-based.
    \item Groups should stay focused on the given prompt during each timed interval.
    \item Written responses should be concise, clear, and connected to lecture concepts.
\end{enumerate} \\
\hline
\end{tabular}
\newpage

\noindent
{\large \textbf{Instructions (detailed activity breakdown below) continued...}} \\
\vspace{0.5cm}
\noindent
\begin{tabular}{|p{0.9\linewidth}|}
\hline
{\large\textbf{Scoring and Feedback}} \\
\begin{enumerate}[label=\alph*.]
    \item Groups earn points for each round based on accuracy, completeness, and explanation quality.
    \item Bonus points may be awarded for:
    \begin{itemize}
        \item Exceptional clarity or creative analogies.
        \item Effective visual representations (diagrams, mechanisms, models).
        \item Strong group participation observed by the instructor.
    \end{itemize}
    \item Post-round debriefs provide immediate correction and conceptual reinforcement.
    \item Optionally maintain a leaderboard to encourage engagement (without high-stakes pressure).
\end{enumerate} \\
\hline
\end{tabular}

\vspace{0.5cm}
\noindent
\begin{tabular}{|p{0.9\linewidth}|}
\hline
{\large\textbf{Wrap-Up and Reflection}} \\
\begin{enumerate}[label=\alph*.]
    \item After all rounds conclude, collect group responses to analyze trends in understanding.
    \item Highlight areas where multiple groups struggled and re-explain key ideas.
    \item Lead a brief reflection discussion addressing:
    \begin{itemize}
        \item Which lecture topics felt most or least comfortable.
        \item How group discussion deepened understanding of the material.
        \item What strategies students plan to use for individual study.
    \end{itemize}
    \item Conclude by connecting the reviewed concepts to upcoming assessments or applications.
\end{enumerate} \\
\hline
\end{tabular}

\newpage
\noindent{\large\textbf{Content (fully worked out problems \& answers)}}\\[0.2em]
\vspace{0.2cm}

\noindent\begin{tabular}{|p{0.9\linewidth}|}
\hline
\textbf{(1)} For the reaction $\text{A} + \text{heat} \rightarrow \text{B}$, which statement is correct? \\[0.2cm]
\begin{minipage}[t]{0.85\linewidth}
\begin{enumerate}[label=\alph*.]
\item Adding heat increases [B]
\item Adding heat decreases [B]
\item Removing heat increases [B]
\item Heat has no effect on equilibrium
\end{enumerate}
\end{minipage}\\[0.2cm]
\hline
\end{tabular}

\noindent\begin{tabular}{|p{0.9\linewidth}|}
\hline
\textbf{(2)} Which of the following best describes a saturated solution at equilibrium?\\[0.2cm]
\begin{minipage}[t]{0.85\linewidth}
\begin{enumerate}[label=\alph*.]
\item Solute is dissolving faster than it crystallizes
\item Solute and solvent are reacting chemically
\item Dissolving and crystallization occur at equal rates
\item No solute particles are moving
\end{enumerate}
\end{minipage}\\[0.2cm]
\hline
\end{tabular}

\noindent\begin{tabular}{|p{0.9\linewidth}|}
\hline
\textbf{(3)} Calcium hydroxide reacts with hydrochloric acid to form calcium chloride and water. Which statement best describes the acid-base behavior occurring in this reaction?\\[0.2cm]
Ca(OH)$_2$ + 2HCl $\rightarrow$ CaCl$_2$ + 2H$_2$O\\[0.2cm]
\begin{minipage}[t]{0.85\linewidth}
\begin{enumerate}[label=\alph*.]
\item HCl donates protons to OH$^-$ ions, forming water as Ca$^{2+}$ pairs with Cl$^-$
\item Ca(OH)$_2$ donates protons to HCl, forming hydrogen gas
\item Both reactants accept protons, preventing water formation
\item Ca$^{2+}$ acts as the proton donor while Cl$^-$ acts as the proton acceptor
\end{enumerate}
\end{minipage}\\[0.2cm]
\hline
\end{tabular}

\noindent\begin{tabular}{|p{0.9\linewidth}|}
\hline
\textbf{(4)} When comparing HF and HCl, HF has a much higher boiling point. Which force best explains this difference?\\[0.2cm]
\begin{minipage}[t]{0.85\linewidth}
\begin{enumerate}[label=\alph*.]
\item Dipole-dipole interactions
\item Hydrogen bonding
\item London dispersion forces
\item Covalent bonding
\end{enumerate}
\end{minipage}\\[0.2cm]
\hline
\end{tabular}

\noindent\begin{tabular}{|p{0.9\linewidth}|}
\hline
\textbf{(5)} Which of the following increases when temperature increases for a gas at constant volume?\\[0.2cm]

\begin{minipage}[t]{0.85\linewidth}
\begin{enumerate}[label=\alph*.]
\item Pressure
\item Molecular mass
\item Intermolecular forces
\item Volume
\end{enumerate}
\end{minipage}\\[0.2cm]
\hline
\end{tabular}

\noindent\begin{tabular}{|p{0.9\linewidth}|}
\hline
\textbf{(6)} A chemist has 50.0~mL of 2.0~M HCl solution. She adds water to dilute it to a final volume of 200.0~mL. What is the molarity of the diluted solution?\\[0.2cm]
\begin{minipage}[t]{0.85\linewidth}
\begin{enumerate}[label=\alph*.]
\item 0.25~M
\item 0.50~M
\item 1.0~M
\item 2.0~M
\end{enumerate}
\end{minipage}\\[0.2cm]
\hline
\end{tabular}

\noindent\begin{tabular}{|p{0.9\linewidth}|}
\hline
\textbf{(7)} An unsaturated solution contains less solute than the maximum that can dissolve at that temperature.\\[0.2cm]
\textbf{True} \hspace{1cm} \textbf{False}\\[0.2cm]
\hline
\end{tabular}

\noindent\begin{tabular}{|p{0.9\linewidth}|}
\hline
\textbf{(8)} What happens to an atom's oxidation number when it is oxidized?\\[0.2cm]
\begin{minipage}[t]{0.85\linewidth}
\begin{enumerate}[label=\alph*.]
\item It increases
\item It decreases
\item It stays the same
\item It becomes zero
\end{enumerate}
\end{minipage}\\[0.2cm]
\hline
\end{tabular}

\noindent\begin{tabular}{|p{0.9\linewidth}|}
\hline
\textbf{(9)} Which of the following solutes is generally soluble in water according to solubility rules?\\[0.2cm]
\begin{minipage}[t]{0.85\linewidth}
\begin{enumerate}[label=\alph*.]
\item BaSO$_4$
\item AgCl
\item KNO$_3$
\item PbI$_2$
\end{enumerate}
\end{minipage}\\[0.2cm]
\hline
\end{tabular}

\noindent\begin{tabular}{|p{0.9\linewidth}|}
\hline
\textbf{(10)} If $Q > K_{eq}$ for the reaction $\text{F} \rightleftharpoons \text{G}$, which direction will the system shift?\\[0.2cm]
\begin{minipage}[t]{0.85\linewidth}
\begin{enumerate}[label=\alph*.]
\item Toward reactants
\item Toward products
\item No shift
\item The reaction stops
\end{enumerate}
\end{minipage}\\[0.2cm]
\hline
\end{tabular}

\noindent\begin{tabular}{|p{0.9\linewidth}|}
\hline
\textbf{(11)} 0.250~mol of oxygen gas is collected at 740.~mmHg and 18.0$^\circ$C in a 2.00~L flask. What would be the pressure at 0.500~mol in same flask and same T? (assume ideal)\\[0.2cm]

\begin{minipage}[t]{0.85\linewidth}
\begin{enumerate}[label=\alph*.]
\item 1480~mmHg
\item 740.~mmHg
\item 370.~mmHg
\item 1110~mmHg
\end{enumerate}
\end{minipage}\\[0.2cm]
\hline
\end{tabular}

\noindent\begin{tabular}{|p{0.9\linewidth}|}
\hline
\textbf{(12)} A Brönsted-Lowry acid is defined as a substance that: \\[0.2cm]
\begin{minipage}[t]{0.85\linewidth}
\begin{enumerate}[label=\alph*.]
\item Donates a proton
\item Accepts a proton
\item Releases electrons
\item Accepts electrons
\end{enumerate}
\end{minipage}\\[0.2cm]
\hline
\end{tabular}

\noindent\begin{tabular}{|p{0.9\linewidth}|}
\hline
\textbf{(13)} A buffer solution is prepared using 0.120~mol NH$_3$ and 0.080~mol NH$_4^+$ with $K_a=5.6\times10^{-10}$ for NH$_4^+$. After calculating the pH, what is the corresponding $[\mathrm{OH^-}]$?[0.2cm]
\begin{minipage}[t]{0.85\linewidth}
\begin{enumerate}[label=\alph*.]
\item pH = 9.32,\ $[\mathrm{OH^-}] = 2.13\times10^{-5}$~M
\item pH = 9.43,\ $[\mathrm{OH^-}] = 1.10\times10^{-5}$~M
\item pH = 9.43,\ $[\mathrm{OH^-}] = 2.68\times10^{-5}$~M
\item pH = 8.90,\ $[\mathrm{OH^-}] = 7.99\times10^{-6}$~M
\end{enumerate}
\end{minipage}[0.2cm]
\hline
\end{tabular}


\noindent\begin{tabular}{|p{0.9\linewidth}|}
\hline
\textbf{(14)} Which of the following would double the (m/v)\% of a solution?\\[0.2cm]
\begin{minipage}[t]{0.85\linewidth}
\begin{enumerate}[label=\alph*.]
\item Double solute, same volume
\item Double volume, same solute
\item Halve solute, double volume
\item Keep solute and volume constant
\end{enumerate}
\end{minipage}\\[0.2cm]
\hline
\end{tabular}

\noindent\begin{tabular}{|p{0.9\linewidth}|}
\hline
\textbf{(15)} A 5\% (m/v) glucose solution contains 5~g of glucose in:\\[0.2cm]
\begin{minipage}[t]{0.85\linewidth}
\begin{enumerate}[label=\alph*.]
\item 5~mL of water
\item 100~mL of water
\item 100~mL of solution
\item 95~mL of solution
\end{enumerate}
\end{minipage}\\[0.2cm]
\hline
\end{tabular}

\noindent\begin{tabular}{|p{0.9\linewidth}|}
\hline
\textbf{(16)} The density of aluminum is $2.70~\text{g/cm}^3$. What is the mass of a block of aluminum with a volume of $15.0~\text{cm}^3$? \\
\begin{minipage}[t]{0.85\linewidth}
\begin{enumerate}[label=\alph*.]
    \item $40.5~\text{g}$
    \item $5.56~\text{g}$
    \item $17.7~\text{g}$
    \item $2.70~\text{g}$
\end{enumerate}
\end{minipage}\\[0.2cm]
\hline
\end{tabular}

\noindent\begin{tabular}{|p{0.9\linewidth}|}
\hline
\textbf{(17)} Convert 0.750 L into cubic centimeters. (Use $1~\text{L} = 1000~\text{cm}^3$) \\
\begin{minipage}[t]{0.85\linewidth}
\begin{enumerate}[label=\alph*.]
    \item $75.0~\text{cm}^3$
    \item $750.~\text{cm}^3$
    \item $7.50~\text{cm}^3$
    \item $0.00750~\text{cm}^3$
\end{enumerate}
\end{minipage}\\[0.2cm]
\hline
\end{tabular}

\noindent\begin{tabular}{|p{0.9\linewidth}|}
\hline
\textbf{(18)} Identify the reaction type:\\[0.2cm]
2H$_2$O$_2$ $\rightarrow$ 2H$_2$O + O$_2$\\[0.2cm]
\begin{minipage}[t]{0.85\linewidth}
\begin{enumerate}[label=\alph*.]
\item Combination
\item Combustion
\item Decomposition
\item Exchange
\end{enumerate}
\end{minipage}\\[0.2cm]
\hline
\end{tabular}

\noindent\begin{tabular}{|p{0.9\linewidth}|}
\hline
\textbf{(19)} What term describes the process by which a solid dissolves until no more can dissolve under given conditions?\\[0.2cm]
\begin{minipage}[t]{0.85\linewidth}
\begin{enumerate}[label=\alph*.]
\item Crystallization
\item Saturation
\item Dilution
\item Precipitation
\end{enumerate}
\end{minipage}\\[0.2cm]
\hline
\end{tabular}

\noindent\begin{tabular}{|p{0.9\linewidth}|}
\hline
\textbf{(20)} A cylinder contains 0.850~mol of gas at 2.50~atm. What is the temperature (K) if the volume is 15.0~L?\\[0.2cm]

\begin{minipage}[t]{0.85\linewidth}
\begin{enumerate}[label=\alph*.]
\item 210.~K
\item 298~K
\item 425~K
\item 537~K
\end{enumerate}
\end{minipage}\\[0.2cm]
\hline
\end{tabular}

\noindent\begin{tabular}{|p{0.9\linewidth}|}
\hline
\textbf{(21)} A student notices that hexane and water do not mix when shaken together. Which type of intermolecular forces explains this behavior?\\[0.2cm]
\begin{minipage}[t]{0.85\linewidth}
\begin{enumerate}[label=\alph*.]
\item London dispersion forces dominate in hexane, but water is polar.
\item Both substances form hydrogen bonds with each other.
\item Hexane is polar, and water is nonpolar.
\item Dipole-dipole interactions attract the two liquids together.
\end{enumerate}
\end{minipage}\\[0.2cm]
\hline
\end{tabular}

\noindent\begin{tabular}{|p{0.9\linewidth}|}
\hline
\textbf{(22)} Which of the following represents a redox reaction?\\[0.2cm]
\begin{minipage}[t]{0.85\linewidth}
\begin{enumerate}[label=\alph*.]
\item HCl + NaOH $\rightarrow$ NaCl + H$_2$O\\[0.1cm]
\item Zn + 2HCl $\rightarrow$ ZnCl$_2$ + H$_2$\\[0.1cm]
\item CaCO$_3$ $\rightarrow$ CaO + CO$_2$\\[0.1cm]
\item NaOH + HNO$_3$ $\rightarrow$ NaNO$_3$ + H$_2$O\\[0.1cm]
\end{enumerate}

\end{minipage}\\[0.2cm]
\hline
\end{tabular}

\noindent\begin{tabular}{|p{0.9\linewidth}|}
\hline
\textbf{(23)} A solution has [H$_3$O$^+$] = $2.22\times10^{-6}$ M. Calculate the pH, pOH, and [OH$^-$]. \\[4cm]
\hline
\end{tabular}

\noindent\begin{tabular}{|p{0.9\linewidth}|}
\hline
\textbf{(24)} A buffer is prepared using 0.250~M HCO$_3^-$ and 0.400~M CO$_3^{2-}$ (p$K_a2 = 10.33$). What are the buffer pH and corresponding pOH?\\[0.2cm]
\begin{minipage}[t]{0.85\linewidth}
\begin{enumerate}[label=\alph*.]
\item pH = 10.13,\ pOH = 3.87
\item pH = 10.33,\ pOH = 3.67
\item pH = 10.73,\ pOH = 3.27
\item pH = 10.15,\ pOH = 3.85
\end{enumerate}
\end{minipage}\\[0.2cm]
\hline
\end{tabular}

\noindent\begin{tabular}{|p{0.9\linewidth}|}
\hline
\textbf{(25)} A 425~mL cylinder contains nitrogen at $2.50\times10^{2}$~kPa and 22.0$^\circ$C. What pressure (in atm) will result if the gas is moved to a 1.00~L container and heated to 125$^\circ$C, assuming moles constant?\\[0.2cm]

\begin{minipage}[t]{0.85\linewidth}
\begin{enumerate}[label=\alph*.]
\item 0.564~atm
\item 1.41~atm
\item 0.282~atm
\item 2.24~atm
\end{enumerate}
\end{minipage}\\[0.2cm]
\hline
\end{tabular}

\noindent\begin{tabular}{|p{0.9\linewidth}|}
\hline
\textbf{(26)} What type of reaction occurs when potassium iodide reacts with lead(II) nitrate to form lead(II) iodide and potassium nitrate?\\[0.2cm]
2KI + Pb(NO$_3$)$_2$ $\rightarrow$ 2KNO$_3$ + PbI$_2$\\[0.2cm]
\begin{minipage}[t]{0.85\linewidth}
\begin{enumerate}[label=\alph*.]
\item Exchange
\item Combination
\item Decomposition
\item Combustion
\end{enumerate}
\end{minipage}\\[0.2cm]
\hline
\end{tabular}

\noindent\begin{tabular}{|p{0.9\linewidth}|}
\hline
\textbf{(27)} When salt dissolves in water, why is this considered a physical change? \\
\begin{minipage}[t]{0.85\linewidth}
\begin{enumerate}[label=\alph*.]
    \item The process creates a new chemical substance with different bonds.
    \item The salt changes state but its chemical composition remains the same.
    \item The salt reacts with water to form a new compound.
    \item The ions rearrange permanently to form a new molecule.
\end{enumerate}
\end{minipage}\\[0.2cm]
\hline
\end{tabular}

\noindent\begin{tabular}{|p{0.9\linewidth}|}
\hline
\textbf{(28)} A buffer is made by mixing 0.200~M HF and 0.350~M NaF (p$K_a=3.17$). After preparing the buffer, what are the pH and resulting $[\mathrm{H_3O^+}]$?\\[0.2cm]
\begin{minipage}[t]{0.85\linewidth}
\begin{enumerate}[label=\alph*.]
\item pH = 2.93,\ $[\mathrm{H_3O^+}] = 1.18\times10^{-3}$~M
\item pH = 3.33,\ $[\mathrm{H_3O^+}] = 4.7\times10^{-4}$~M
\item pH = 3.17,\ $[\mathrm{H_3O^+}] = 6.8\times10^{-4}$~M
\item pH = 3.60,\ $[\mathrm{H_3O^+}] = 2.5\times10^{-4}$~M
\end{enumerate}
\end{minipage}\\[0.2cm]
\hline
\end{tabular}

\noindent\begin{tabular}{|p{0.9\linewidth}|}
\hline
\textbf{(29)} A solution has [$H_3O^+$] = $3.20\times10^{-6}$ M. Its pH is closest to: \\[0.2cm]
\begin{minipage}{0.85\linewidth}
\begin{enumerate}[label=\alph*.]
\item 4.49 \item 5.49 \item 6.49 \item 7.49
\end{enumerate}
\end{minipage}\\[0.2cm]
\hline
\end{tabular}

\noindent\begin{tabular}{|p{0.9\linewidth}|}
\hline
\textbf{(30)} A weak acid has Ka = $3.00\times10^{-7}$. Which pKa is closest? \\[0.2cm]
\begin{minipage}{0.85\linewidth}
\begin{enumerate}[label=\alph*.]
\item 3.96 \item 4.22 \item 5.44 \item 6.80
\end{enumerate}
\end{minipage}\\[0.2cm]
\hline
\end{tabular}

\noindent\begin{tabular}{|p{0.9\linewidth}|}
\hline
\textbf{(31)} Two rigid containers are connected by a closed valve. One container holds helium at 2.0~atm in 4.0~L, and the other holds neon at 4.0~atm in 2.0~L. When the valve is opened and gases mix, what is the total pressure in the combined 6.0~L system (assume constant temperature)?\\[0.2cm]
\begin{minipage}[t]{0.85\linewidth}
\begin{enumerate}[label=\alph*.]
\item 2.0~atm
\item 2.7~atm
\item 4.0~atm
\item 6.0~atm
\end{enumerate}
\end{minipage}\\[0.2cm]
\hline
\end{tabular}

\noindent\begin{tabular}{|p{0.9\linewidth}|}
\hline
\textbf{(32)} A solution with pH = 2.00 has a hydronium concentration that is how many times larger than a solution with pH = 4.00? \\[0.2cm]
\begin{minipage}{0.85\linewidth}
\begin{enumerate}[label=\alph*.]
\item 2x \item 10x \item 100x \item 1000x
\end{enumerate}
\end{minipage}\\[0.2cm]
\hline
\end{tabular}

\noindent\begin{tabular}{|p{0.9\linewidth}|}
\hline
\textbf{(33)} Which of the following reactions is a redox reaction?\\[0.2cm]
\begin{minipage}[t]{0.85\linewidth}
\begin{enumerate}[label=\alph*.]
\item Zn + CuSO$_4$ $\rightarrow$ ZnSO$_4$ + Cu\\[0.1cm]
\item AgNO$_3$ + NaCl $\rightarrow$ AgCl + NaNO$_3$\\[0.1cm]
\item CaCO$_3$ $\rightarrow$ CaO + CO$_2$\\[0.1cm]
\item H$_2$O + SO$_3$ $\rightarrow$ H$_2$SO$_4$\\[0.1cm]
\end{enumerate}
\end{minipage}\\[0.2cm]
\hline
\end{tabular}

\noindent\begin{tabular}{|p{0.9\linewidth}|}
\hline
\textbf{(34)} Which factor determines how much the freezing point of a solution is lowered?\\[0.2cm]
\begin{minipage}[t]{0.85\linewidth}
\begin{enumerate}[label=\alph*.]
\item The number of solute particles
\item The solute's color
\item The solvent's taste
\item The solute's shape
\end{enumerate}
\end{minipage}\\[0.2cm]
\hline
\end{tabular}

\noindent\begin{tabular}{|p{0.9\linewidth}|}
\hline
\textbf{(35)} What happens to the molarity of a solution when it is diluted with more solvent?\\[0.2cm]
\begin{minipage}[t]{0.85\linewidth}
\begin{enumerate}[label=\alph*.]
\item It increases
\item It decreases
\item It stays the same
\item It doubles
\end{enumerate}
\end{minipage}\\[0.2cm]
\hline
\end{tabular}

\noindent\begin{tabular}{|p{0.9\linewidth}|}
\hline
\textbf{(36)} What is the molarity of a solution containing 5.00~mol of KBr in 2.00~L of solution?\\[0.2cm]
\begin{minipage}[t]{0.85\linewidth}
\begin{enumerate}[label=\alph*.]
\item 2.00~M
\item 2.50~M
\item 5.00~M
\item 10.0~M
\end{enumerate}
\end{minipage}\\[0.2cm]
\hline
\end{tabular}

\noindent\begin{tabular}{|p{0.9\linewidth}|}
\hline
\textbf{(37)} What is the percent yield if 10.0~g of product are expected but only 8.20~g are obtained?\\[0.2cm]

\begin{minipage}[t]{0.85\linewidth}
\begin{enumerate}[label=\alph*.]
\item 82.0\%
\item 92.0\%
\item 80.0\%
\item 88.0\%
\end{enumerate}
\end{minipage}\\[0.2cm]
\hline
\end{tabular}

\noindent\begin{tabular}{|p{0.9\linewidth}|}
\hline
\textbf{(38)} Which species is the conjugate acid of NH$_3$? \\[0.2cm]
\begin{minipage}{0.85\linewidth}
\begin{enumerate}[label=\alph*.]
\item NH$_2^-$ \item NH$_4^+$ \item N$^{3-}$ \item NH$_2$OH
\end{enumerate}
\end{minipage}\\[0.2cm]
\hline
\end{tabular}

\noindent\begin{tabular}{|p{0.9\linewidth}|}
\hline
\textbf{(39)} Osmosis is the movement of solvent molecules from:\\[0.2cm]
\begin{minipage}[t]{0.85\linewidth}
\begin{enumerate}[label=\alph*.]
\item Low solute concentration to high solute concentration
\item High solute concentration to low solute concentration
\item High pressure to low pressure
\item High temperature to low temperature
\end{enumerate}
\end{minipage}\\[0.2cm]
\hline
\end{tabular}

\noindent\begin{tabular}{|p{0.9\linewidth}|}
\hline
\textbf{(40)} A balloon is filled with a gas mixture containing 0.6~mol O$_2$ and 0.4~mol N$_2$ at a total pressure of 1.0~atm. What is the partial pressure of O$_2$?\\[0.2cm]
\begin{minipage}[t]{0.85\linewidth}
\begin{enumerate}[label=\alph*.]
\item 0.4~atm
\item 0.6~atm
\item 1.0~atm
\item 1.5~atm
\end{enumerate}
\end{minipage}\\[0.2cm]
\hline
\end{tabular}

\noindent\begin{tabular}{|p{0.9\linewidth}|}
\hline
\textbf{(41)} How many grams of KNO$_3$ are required to prepare 250.0~mL of 0.200~M KNO$_3$ solution?\\[0.2cm]
\begin{minipage}[t]{0.85\linewidth}
\begin{enumerate}[label=\alph*.]
\item 2.53~g
\item 5.06~g
\item 10.1~g
\item 20.2~g
\end{enumerate}
\end{minipage}\\[0.2cm]
\hline
\end{tabular}

\noindent\begin{tabular}{|p{0.9\linewidth}|}
\hline
\textbf{(42)} Convert 250.0 cm/s to mi/h. (Use $1~\text{mi}=1.609~\text{km}$) \\
\begin{minipage}[t]{0.85\linewidth}
\begin{enumerate}[label=\alph*.]
    \item $5.59~\text{mi/h}$
    \item $15.6~\text{mi/h}$
    \item $0.155~\text{mi/h}$
    \item $1.56~\text{mi/h}$
\end{enumerate}
\end{minipage}\\[0.2cm]
\hline
\end{tabular}

\noindent\begin{tabular}{|p{0.9\linewidth}|}
\hline
\textbf{(43)} For the reaction $\text{R} + \text{S} \rightleftharpoons \text{RS}$, which change could temporarily increase Q?\\[0.2cm]
\begin{minipage}[t]{0.85\linewidth}
\begin{enumerate}[label=\alph*.]
\item Adding RS
\item Removing R
\item Removing S
\item Diluting all species equally
\end{enumerate}
\end{minipage}\\[0.2cm]
\hline
\end{tabular}

\noindent\begin{tabular}{|p{0.9\linewidth}|}
\hline
\textbf{(44)} In the reaction $\text{HCl} + \text{H}_2\text{O} \rightarrow \text{H}_3\text{O}^+ + \text{Cl}^-$, which is the Brönsted-Lowry acid?\\[0.2cm]
\begin{minipage}[t]{0.85\linewidth}
\begin{enumerate}[label=\alph*.]
\item HCl
\item H$_2$O
\item H$_3$O$^+$
\item Cl$^-$
\end{enumerate}
\end{minipage}\\[0.2cm]
\hline
\end{tabular}

\noindent\begin{tabular}{|p{0.9\linewidth}|}
\hline
\textbf{(45)} Which of the following changes how much a solution’s freezing point is lowered?\\[0.2cm]
\begin{minipage}[t]{0.85\linewidth}
\begin{enumerate}[label=\alph*.]
\item How much solute is added
\item The type of solute (number of particles it makes)
\item A \& B
\item The color of the solute
\end{enumerate}
\end{minipage}\\[0.2cm]
\hline
\end{tabular}

\noindent\begin{tabular}{|p{0.9\linewidth}|}
\hline
\textbf{(46)} Convert $745~\text{torr}$ to atm. (Use $1~\text{atm}=760~\text{torr}$) \\
\begin{minipage}[t]{0.85\linewidth}
\begin{enumerate}[label=\alph*.]
    \item $0.980~\text{atm}$
    \item $1.02~\text{atm}$
    \item $0.760~\text{atm}$
    \item $0.745~\text{atm}$
\end{enumerate}
\end{minipage}\\[0.2cm]
\hline
\end{tabular}

\noindent\begin{tabular}{|p{0.9\linewidth}|}
\hline
\textbf{(47)} A 25\% (v/v) ethanol solution means:\\[0.2cm]
\begin{minipage}[t]{0.85\linewidth}
\begin{enumerate}[label=\alph*.]
\item 25~mL ethanol in 100~mL solution
\item 25~mL ethanol in 100~mL solvent
\item 25~g ethanol in 100~mL solution
\item 25~mL solution in 100~mL ethanol
\end{enumerate}
\end{minipage}\\[0.2cm]
\hline
\end{tabular}

\noindent\begin{tabular}{|p{0.9\linewidth}|}
\hline
\textbf{(48)} Which species is amphiprotic? \\[0.2cm]
\begin{minipage}{0.85\linewidth}
\begin{enumerate}[label=\alph*.]
\item OH$^-$ \item HSO$_4^-$ \item Cl$^-$ \item NH$_4^+$
\end{enumerate}
\end{minipage}\\[0.2cm]
\hline
\end{tabular}

\noindent\begin{tabular}{|p{0.9\linewidth}|}
\hline
\textbf{(49)} A buffer is prepared by mixing 0.300~mol acetic acid (p$K_a=4.76$) with 0.450~mol sodium acetate in 1.00~L solution. What is the pH of the buffer, and what is the corresponding pOH?\\[0.2cm]
\begin{minipage}[t]{0.85\linewidth}
\begin{enumerate}[label=\alph*.]
\item pH = 4.94,\ pOH = 9.06
\item pH = 4.76,\ pOH = 9.24
\item pH = 4.58,\ pOH = 9.42
\item pH = 5.12,\ pOH = 8.88
\end{enumerate}
\end{minipage}\\[0.2cm]
\hline
\end{tabular}

\noindent\begin{tabular}{|p{0.9\linewidth}|}
\hline
\textbf{(50)} For the equilibrium $ \text{PCl}_5(g) \rightleftharpoons \text{PCl}_3(g) + \text{Cl}_2(g)$, increasing pressure causes the reaction to shift: \\[0.2cm]
\begin{minipage}[t]{0.85\linewidth}
\begin{enumerate}[label=\alph*.]
\item Left, toward PCl$_5$
\item Right, toward PCl$_3$ and Cl$_2$
\item No shift because gases are unaffected by pressure
\item Toward whichever side has the higher total mass
\end{enumerate}
\end{minipage}\\[0.2cm]
\hline
\end{tabular}


\noindent\begin{tabular}{|p{0.9\linewidth}|}
\hline
\textbf{(51)} For each of the following, select the option that correctly identifies each property as physical (P) or chemical (C): \
\begin{minipage}[t]{0.85\linewidth}
\begin{enumerate}[label=\alph*.]
\item melting point of ice
\item flammability of hydrogen gas
\item density of aluminum
\item ability of iron to rust
\end{enumerate}
\textbf{Choices:} \
A) P, C, P, C \
B) C, P, C, P \
C) P, P, C, C \
D) C, C, P, P
\end{minipage}[0.2cm]
\hline
\end{tabular}

\noindent\begin{tabular}{|p{0.9\linewidth}|}
\hline
\textbf{(52)} For each of the following, select the option that correctly classifies each as a pure substance (PS) or a mixture (M): \
\begin{minipage}[t]{0.85\linewidth}
\begin{enumerate}[label=\alph*.]
\item sodium chloride
\item air
\item copper
\item seawater
\end{enumerate}
\textbf{Choices:} \
A) PS, M, PS, M \
B) M, PS, M, PS \
C) PS, PS, M, M \
D) M, M, PS, PS
\end{minipage}[0.2cm]
\hline
\end{tabular}

\noindent\begin{tabular}{|p{0.9\linewidth}|}
\hline
\textbf{(53)} A gas occupies 500.~mL at 25$^\circ$C. What volume will it occupy at 100.$^\circ$C if the pressure remains constant? \\[0.2cm]

\begin{minipage}[t]{0.85\linewidth}
\begin{enumerate}[label=\alph*.]
\item 583~mL
\item 625~mL
\item 665~mL
\item 740.~mL
\end{enumerate}
\end{minipage}\\[0.2cm]
\hline
\end{tabular}

\noindent\begin{tabular}{|p{0.9\linewidth}|}
\hline
\textbf{(54)} Convert 5850~$\text{cm}^2$ to $\text{m}^2$. \\
\begin{minipage}[t]{0.85\linewidth}
\begin{enumerate}[label=\alph*.]
    \item $0.0585~\text{m}^2$
    \item $5.85~\text{m}^2$
    \item $585~\text{m}^2$
    \item $0.000585~\text{m}^2$
\end{enumerate}
\end{minipage}\\[0.2cm]
\hline
\end{tabular}

\noindent\begin{tabular}{|p{0.9\linewidth}|}
\hline
\textbf{(55)} Which statement correctly defines ionization for a base in water? \\[0.2cm]
\begin{minipage}[t]{0.85\linewidth}
\begin{enumerate}[label=\alph*.]
\item Forming OH$^-$ by reacting chemically with water rather than releasing pre-existing ions
\item Separating into pre-existing ions without chemical change
\item Producing H$_3$O$^+$ by transferring a proton to water
\item Splitting into neutral molecules without forming ions
\end{enumerate}
\end{minipage}\\[0.2cm]
\hline
\end{tabular}

\noindent\begin{tabular}{|p{0.9\linewidth}|}
\hline
\textbf{(56)} Which term describes an acid that ionizes partially in water? \\[0.2cm]
\begin{minipage}[t]{0.85\linewidth}
\begin{enumerate}[label=\alph*.]
\item Strong acid \item Weak acid \item Strong base \item Neutral compound
\end{enumerate}
\end{minipage}\\[0.2cm]
\hline
\end{tabular}

\noindent\begin{tabular}{|p{0.9\linewidth}|}
\hline
\textbf{(57)} Osmotic pressure depends on which of the following factors?\\[0.2cm]
\begin{minipage}[t]{0.85\linewidth}
\begin{enumerate}[label=\alph*.]
\item The number of solute particles in solution
\item The color of the solute
\item The shape of the solvent molecules
\item The size of the container
\end{enumerate}
\end{minipage}\\[0.2cm]
\hline
\end{tabular}

\noindent\begin{tabular}{|p{0.9\linewidth}|}
\hline
\textbf{(58)} A buffer is formed by mixing 0.500~mol formic acid (p$K_a=3.75$) with 0.300~mol sodium formate. What are the buffer pH and the resulting $[\mathrm{H_3O^+}]$?\\[0.2cm]
\begin{minipage}[t]{0.85\linewidth}
\begin{enumerate}[label=\alph*.]
\item pH = 3.52,\ $[\mathrm{H_3O^+}] = 3.0\times10^{-4}$~M
\item pH = 3.75,\ $[\mathrm{H_3O^+}] = 1.8\times10^{-4}$~M
\item pH = 3.32,\ $[\mathrm{H_3O^+}] = 4.8\times10^{-4}$~M
\item pH = 3.90,\ $[\mathrm{H_3O^+}] = 1.3\times10^{-4}$~M
\end{enumerate}
\end{minipage}\\[0.2cm]
\hline
\end{tabular}

\noindent\begin{tabular}{|p{0.9\linewidth}|}
\hline
\textbf{(59)} A tectonic plate moves 860 km per day. What is its speed in m/s? \\
\begin{minipage}[t]{0.85\linewidth}
\begin{enumerate}[label=\alph*.]
    \item $9.95~\text{m/s}$
    \item $0.0995~\text{m/s}$
    \item $9.95 \times 10^3~\text{m/s}$
    \item $2.47~\text{m/s}$
\end{enumerate}
\end{minipage}\\[0.2cm]
\hline
\end{tabular}

\noindent\begin{tabular}{|p{0.9\linewidth}|}
\hline
\textbf{(60)} Which salt will produce a basic solution upon dissolution due to its parent species? \\[0.2cm]
\begin{minipage}{0.85\linewidth}
\begin{enumerate}[label=\alph*.]
\item NaF \item NH$_4$Br \item KCl \item Ca(NO$_3$)$_2$
\end{enumerate}
\end{minipage}\\[0.2cm]
\hline
\end{tabular}


\newpage
\noindent
{\large\textbf{Checking for Understanding (questions \& answers)}}\\

\begin{tabular}{|p{0.9\linewidth}|}
\hline
\begin{itemize}
    \item How can molarity be used to determine the number of moles of solute present in a given volume of solution?
    \item What steps are required to calculate mass percent, volume percent, or mass/volume percent from experimental data?
    \item How do dilution calculations relate initial and final concentrations and volumes of a solution?
    \item What evidence indicates that a solution is saturated, unsaturated, or supersaturated at a particular temperature?
    \item How do solubility rules help predict whether a precipitate will form during a double replacement reaction?
    \item Why does increasing temperature typically increase the solubility of most solid solutes but decrease the solubility of gases?
    \item How does solute–solvent polarity affect dissolution and miscibility?
    \item How does adding a nonvolatile solute affect the freezing point of a solvent?
    \item Why does dissolving an ionic compound typically lower the freezing point more than dissolving a covalent compound?
    \item What factors determine the magnitude of freezing point depression in a solution?
\end{itemize}\\[0.2cm]
\end{tabular}
\newpage
\begin{tabular}{|p{0.9\linewidth}|}
\hline
\begin{itemize}
    \item Why does the identity of a solute, not just its amount, influence colligative properties?
    \item How does the boiling point of a solution compare to the boiling point of its pure solvent?
    \item Why does adding solute increase the boiling point of a liquid?
    \item Which physical property remains unchanged when the freezing point of a solution is lowered?
    \item Why does temperature itself not change the amount of freezing point depression?
    \item Why do electrolytes and nonelectrolytes produce different magnitudes of colligative effects?
    \item How can a balanced chemical equation be used to determine mole ratios for stoichiometric calculations?
    \item Why is balancing chemical equations required to satisfy conservation of mass?
    \item How can stoichiometric relationships be used to determine theoretical yield?
    \item How can you identify the limiting reactant when given initial quantities of multiple reactants?
    \item Why do ionic compounds dissociate into ions in water, while many molecular compounds do not?
    \item How does solubility influence whether certain ions in a double replacement reaction form a solid precipitate?
    \item Why is water considered an effective solvent for ionic and polar solutes?
    \item How does molecular polarity explain why nonpolar substances like oil do not mix with water?
    \item Why do stronger intermolecular forces correspond to higher boiling points?
\end{itemize}\\[0.2cm]
\end{tabular}
\newpage
\begin{tabular}{|p{0.9\linewidth}|}
\hline
\begin{itemize}
    \item How can temperature changes be used to determine whether a process is endothermic or exothermic?
    \item Why is combustion classified as an exothermic oxidation reaction?
    \item What evidence suggests a chemical change has occurred rather than a physical change?
    \item Why is neutralization considered a double replacement reaction, and how does it differ from precipitation and redox reactions?
    \item How can synthesis, decomposition, single replacement, and double replacement reactions be distinguished based on reactants and products?
    \item How does the activity series help predict whether a single replacement reaction will occur?
    \item How do gas laws relate pressure, volume, temperature, and number of moles for a gas sample?
    \item How does the kinetic molecular theory explain changes in gas pressure and temperature?
\end{itemize}\\[0.2cm]
\end{tabular}
\newpage
\begin{tabular}{|p{0.9\linewidth}|}
\hline
\begin{itemize}
    \item Why must unit conversions and significant-figure rules be followed carefully in multistep calculations?
    \item How do fixed instrumental uncertainties influence how measurements should be recorded?
    \item How can Avogadro’s number be applied to convert between moles and number of particles?
    \item How do electron orbitals describe the probable locations of electrons within atoms?
    \item What is the difference between dissociation and ionization in aqueous solutions?
    \item How can pH, pOH, [H$_3$O$^+$], and [OH$^-$] be calculated from one another?
    \item How can percent ionization be determined for a weak acid given its concentration and $K_a$?
    \item How do $K_a$ and p$K_a$ values indicate the strength of an acid?
    \item How do conjugate acid–base pairs differ within a Brønsted–Lowry framework?
    \item How do polyprotic acids undergo multiple ionization steps, and how does this affect solution pH?
    \item How can salt hydrolysis be used to predict whether a solution formed from a dissolved salt will be acidic, basic, or neutral?
    \item How can consistency checks and correct units ensure that quantitative solution calculations are valid?
\end{itemize} \\[0.2cm]
\hline
\end{tabular}



\newpage

\section*{\underline{Activity 2}}
\vspace{0.2cm}
\begin{tabular}{|p{4.2cm}|p{9.8cm}|}
\hline
\textbf{Collaborative Learning Techniques} & Jigsaw\\
\hline
\textbf{Learning Strategy} & Predict Test Questions \\
\hline
\textbf{Time} & 12:05---12:30\\
\hline
\textbf{Materials} & Pencil, eraser, calculator, periodic table \\
\hline
\textbf{Lecture/Textbook/Old Plans/Slide References} & Chapter 10: Acid and Bases \\
\hline
\end{tabular}
\vspace{0.5cm}

{\large\textbf{Definition}}\\
\noindent
\begin{tabular}{|p{0.9\linewidth}|}
\hline
\begin{enumerate}[label=\alph*.]
    \item\textbf{Jigsaw:} Jigsaws, when used properly, make the group as a whole dependent upon all of them in subgroups. Each group provides a peice of the puzzle. Group members are broken into smaller groups. Each small group works on some aspect of the same problem, question, or issue. They then share their part of the puzzle with the large group.
    \item\textbf{Predict Test Questions:}  Put students in groups of 2 or 3 and assign them to write a test question for a specific topic, ensuring that all topics have been convered. Ask students to write their questions on the board or on an overhead for discussion (would the professor ask the question? What is the answer? etc.). Students will have the benefit of learning to think like the teacher and they'll be able to see additional questions that other students have writen.
 
\end{enumerate} \\
\hline
\end{tabular}
\newpage

\vspace{0.5cm}
\noindent
{\large \textbf{Instructions (detailed activity breakdown below)}} \\

\begin{tabular}{|p{0.9\linewidth}|}
\hline
\begin{enumerate}[label=\alph*.]
    \item \textbf{Setup:} 
    \begin{itemize}
        \item Divide students into small groups and assign each group a role structure (e.g., scribe, presenter, checker, timekeeper).
        \item Provide each group with a blank question card for writing their original test question.
        \item Give each student an individual worksheet to work on while the instructor solves the group-generated questions.
    \end{itemize}

    \item \textbf{Formulating the Question (5 minutes):}
    \begin{itemize}
        \item Each group collaborates to write one test question related to current unit topics (calculations, concepts, reaction types, etc.).
        \item The question must include:
            \begin{itemize}
                \item A clearly stated prompt.
                \item Any necessary data or diagrams.
                \item A final answer key written \emph{lightly} on the back for instructor use only.
            \end{itemize}
        \item Groups must agree unanimously that the question is fair, solvable, and aligned with recent course content.
        \item When time is up, the instructor collects all question cards and begins solving them.
    \end{itemize}

    \item \textbf{Independent Work Period (Instructor Solves Questions):}
    \begin{itemize}
        \item While the instructor works through each group’s question, students complete their individual worksheets quietly.
        \item Groups may discuss the worksheet only after all members finish the same section.
        \item The instructor prepares corrected solutions and confirms the correct answer for each group-generated test question.
    \end{itemize}
\end{enumerate}\\[0.2cm]
\end{tabular}
\newpage
\begin{tabular}{|p{0.9\linewidth}|}
\hline
\begin{enumerate}
    \item \textbf{Question Exchange and Group Challenge:}
    \begin{itemize}
        \item Once solved, the instructor returns each group \emph{another group’s} test question.
        \item Groups must work collaboratively using assigned roles to solve the received question.
        \item All members must agree on the final answer before reporting it.
        \item Each group provides its final answer verbally to the instructor when ready.
    \end{itemize}

    \item \textbf{Scoring:}
    \begin{itemize}
        \item A group earns \textbf{1 point} for submitting the correct answer to the exchanged question.
        \item A group earns \textbf{1 bonus point} if their originally written test question is solved correctly by the receiving group.
        \item Groups with full agreement and strong justification may receive optional teamwork bonus credit.
        \item The team with the highest point total wins the challenge.
    \end{itemize}

    \item \textbf{Rules:}
    \begin{itemize}
        \item Groups must produce their question collaboratively; no individual writing without group input.
        \item All group members must agree on submitted answers—consensus is required.
        \item Worksheet work during the instructor-solving phase must be independent.
        \item Direct sharing of answers between groups is prohibited.
        \item Questions must be original and may not copy textbook or online sources verbatim.
    \end{itemize}

    \item \textbf{Wrap-Up:}
    \begin{itemize}
        \item Review a selection of the student-generated questions as a whole class.
        \item Highlight creative question design, common misconceptions, and strong reasoning.
        \item Discuss how generating questions reinforced deeper understanding of unit concepts.
        \item Encourage reflection on teamwork, consensus-building, and accuracy in problem solving.
    \end{itemize}
\end{enumerate} \\
\hline
\end{tabular}

\\
\hline
\end{tabular}
\newpage
\noindent{\large\textbf{Content (fully worked out problems \& answers)}}\\[0.2em]
\vspace{0.2cm}

\begin{tabular}{|p{0.9\linewidth}|}
\hline
\begin{enumerate}[label=\arabic*.]

    \item \textbf{Chapter 1 – Basic Concepts About Matter}
    \begin{itemize}
        \item Classification of matter: element, compound, mixture (simple flow diagram).
        \item States of matter and particle spacing.
        \item Physical vs. chemical properties and changes (icon examples).
        \item Density formula and basic interpretation (mass/volume block sketch).
    \end{itemize}

    \item \textbf{Chapter 2 – Measurements in Chemistry}
    \begin{itemize}
        \item Metric units used in chemistry (g, L, m) and common prefixes (k, c, m).
        \item Significant figure rules (quick chart: zeros counting vs. not).
        \item Dimensional analysis (one-step unit conversion map).
        \item Temperature scales: Celsius and Kelvin (short conversion diagram).
    \end{itemize}

    \item \textbf{Chapter 3 – Atomic Structure and the Periodic Table}
    \begin{itemize}
        \item Structure of the atom: protons, neutrons, electrons (basic Bohr-style diagram).
        \item Atomic number, mass number, and isotopes (notation with symbols).
        \item Organization of periodic table: groups, periods, metals/nonmetals.
        \item Very basic periodic patterns: valence electrons by group, simple charge patterns for ions.
    \end{itemize}

    \item \textbf{Chapter 4 – Ionic Bonding}
    \begin{itemize}
        \item How ions form (metal→cation, nonmetal→anion).
        \item Simple electron transfer diagrams.
        \item Writing ionic formulas using charge balance.
        \item Naming ionic compounds (metal + nonmetal ending in -ide).
        \item Recognizing common polyatomic ions (small table only).
    \end{itemize}

    \item \textbf{Chapter 5 – Covalent Bonding}
    \begin{itemize}
        \item Covalent bonds as shared electrons (simple dot diagrams).
        \item Single vs. double bonds (basic Lewis examples).
        \item Electronegativity and polar covalent bonds (arrow-with-δ symbols).
        \item VSEPR shapes: linear, bent, tetrahedral 
        \item Naming simple molecular compounds (prefix system).
    \end{itemize}
\end{enumerate}\\[0.2cm]
\end{tabular}
\newpage
\begin{tabular}{|p{0.9\linewidth}|}
\hline
\begin{enumerate}
    \item \textbf{Chapter 6 – Chemical Calculations}
    \begin{itemize}
        \item Formula mass / molar mass (periodic table extraction).
        \item The mole and Avogadro’s number (conceptual diagram with “counting unit”).
        \item Mole conversions (grams ↔ moles ↔ particles).
        \item Balanced equations and mole ratios (visual ratio arrows).
        \item Basic limiting reactant setups (simple table comparing amounts).
    \end{itemize}

    \item \textbf{Chapter 7 – Gases, Liquids, and Solids}
    \begin{itemize}
        \item Particle diagrams for each state of matter.
        \item Gas variables (P, V, T, n) and direct/inverse relationships.
        \item Basic gas laws: Boyle’s, Charles’, Combined Gas Law (no derivations).
        \item Intermolecular forces: dispersion, dipole, H-bonding (simple ranking diagram).
        \item Heating curve sketch (solid→liquid→gas steps).
    \end{itemize}

    \item \textbf{Chapter 8 – Solutions}
    \begin{itemize}
        \item Solute vs. solvent (particle picture).
        \item Types of solutions: electrolyte vs. nonelectrolyte (simple ion break-apart diagram).
        \item Concentrations: molarity, mass\%, volume\%, mass/volume\%.
        \item Dilutions (M₁V₁ = M₂V₂).
        \item Solubility and saturation (curve showing solubility vs. temperature).
        \item Colligative properties at an intro level: FPD and BPE (directional arrows only).
    \end{itemize}

    \item \textbf{Chapter 9 – Chemical Reactions}
    \begin{itemize}
        \item Evidence of a reaction (gas, color change, precipitate, energy change).
        \item Reaction types: synthesis, decomposition, single replacement, double replacement, combustion.
        \item Predicting products using simple patterns (A + B → AB, AB → A + B, etc.).
        \item Precipitation diagram (two aqueous ions forming a solid).
        \item Energy diagrams: exothermic vs. endothermic (simple up/down curve).
    \end{itemize}

    \item \textbf{Chapter 10 – Acids, Bases, and Salts}
    \begin{itemize}
        \item Properties of acids and bases (taste/feel cautions + conductivity).
        \item Arrhenius definitions (H⁺ producers, OH⁻ producers).
        \item Very basic Brønsted concept (acid donates H⁺).
        \item pH scale (0–14 color bar).
        \item Neutralization reaction (acid + base → salt + water).
        \item Strong vs. weak acids (simple dissociation comparison).
        \item Salt solutions: acidic, basic, neutral (simple prediction table).
    \end{itemize}

\end{enumerate} \\
\hline
\end{tabular}


\newpage
\noindent
{\large\textbf{Checking for Understanding (questions)}}\\[0.2em]
\noindent
\begin{tabular}{|p{0.9\linewidth}|}
\hline
\begin{enumerate}[label=\alph*.]

\item How do you distinguish between elements, compounds, and mixtures using particle-level representations?
\item How can you determine whether a change is physical or chemical based on observable evidence?
\item How do you calculate density from mass and volume, and how can density be used to identify an unknown substance?
\item How do metric prefixes relate to base units, and how do you convert between units using factor–label (dimensional analysis) methods?
\item How do you determine the number of significant figures in a measurement, and how do rounding rules differ for multiplication/division versus addition/subtraction?
\item How do you correctly record the value on a graduated cylinder, thermometer, or balance, including the estimated digit?
\item How does the atomic number determine an element’s identity, and how do isotopes differ in atomic mass but not chemical behavior?
\item How do you calculate the number of protons, neutrons, and electrons in neutral atoms and simple ions?
\item How does the arrangement of the periodic table help predict the charges of ions formed by main-group elements?
\item How do metals and nonmetals form ionic compounds through electron transfer, and how do you write formulas that balance charges?
\item How do you draw simple Lewis dot structures for main-group elements and use them to predict bonding patterns?
\item How do differences in electronegativity determine whether bonds are nonpolar covalent, polar covalent, or ionic?
\item How do molecular shapes such as linear, bent, trigonal pyramidal, and tetrahedral influence molecular polarity?
\item How do you calculate molar mass from a chemical formula, and how do you use it to convert between grams, moles, and number of particles?
\item How do you balance a chemical equation, and why is balancing necessary to satisfy the law of conservation of mass?
\item How do you use mole ratios from a balanced equation to relate quantities of reactants and products?
\item How do you determine the limiting reactant by comparing available moles of each reactant, and how do you use it to calculate theoretical yield?
\end{enumerate}\\[0.2cm]
\end{tabular}
\newpage
\begin{tabular}{|p{0.9\linewidth}|}
\hline
\begin{enumerate}
\item How do particle diagrams illustrate the differences between gases, liquids, and solids in terms of spacing and motion?
\item How do gas laws such as Boyle’s and Charles’s describe relationships between P, V, and T, and how do you choose the correct law for a given situation?
\item How does increased kinetic energy at higher temperatures influence gas pressure, diffusion, and average molecular speed?
\item How do you distinguish electrolytes from nonelectrolytes by their behavior in water, and how do ion diagrams represent dissociation?
\item How do you calculate concentration using molarity, mass percent, volume percent, and mass/volume percent for different types of solutions?
\item How do you determine whether a solution is unsaturated, saturated, or supersaturated using a solubility curve?
\item How do you apply solubility rules to predict whether a precipitate forms when two aqueous solutions are mixed?
\item How do you compute dilution values using $M_{1}V_{1} = M_{2}V_{2}$, and what quantity remains constant during a dilution process?
\item How does the number of dissolved particles affect boiling point elevation and freezing point depression, and why do ionic compounds cause larger changes?
\item How can you classify reactions as synthesis, decomposition, single replacement, double replacement, or combustion based on their reactants and products?
\item How do you identify whether energy is absorbed or released by examining simple reaction diagrams for endothermic and exothermic changes?
\item How do Arrhenius acids and bases differ, and how does the formation of hydronium allow you to represent acids in aqueous solution?
\item How do you determine pH from hydronium concentration, and how do you convert between pH, pOH, $[H_{3}O^{+}]$, and $[OH^{-}]$ using $K_w$?
\item How do strong acids and strong bases behave differently from weak ones in terms of dissociation and resulting pH?
\item How do you identify conjugate acid–base pairs in a Brønsted–Lowry reaction, and how do proton transfers define these relationships?
\item How do diprotic and triprotic acids ionize in steps, and why does each successive step release fewer hydronium ions?
\item How do salts formed from strong or weak acids and bases influence the pH of a solution through hydrolysis?
\item How do you verify that calculated pH, pOH, $[H_{3}O^{+}]$, and $[OH^{-}]$ are internally consistent and properly labeled with units or unitless conventions?

\end{enumerate}
\hline  
\end{tabular}

\newpage
\noindent
\section*{\underline{Closer}}
\begin{tabular}{|p{4.2cm}|p{9.8cm}|}
\hline
\textbf{Collaborative Learning Techniques} & Group Survey\\
\hline
\textbf{Learning Strategy} & Brain dump\\
\hline
\textbf{Time} & 12:30---12:32\\
\hline
\textbf{Materials} & Index cards, writing utensils\\
\hline
\textbf{Lecture/Textbook/Old Plans/Slide References} & Chapters 1-10 \\
\hline
\end{tabular}

\vspace{0.5cm}
\noindent
\begin{tabular}{|p{0.9\linewidth}|}
\hline
{\large\textbf{Definitions}}
\begin{enumerate}[label=\alph*.]
   \item\textbf{Group survey:} Each group member is surveyed to discover their position on an issue, problem or topic. This process ensures that each member of the group is allowed to offer or state their point of view.
   \item\textbf{Brain dump:} This is a great closer for students to sum up their knowledge on the topic(s) from the session. Have the students take out a blank piece of paper /Google Doc and write/type("dump") everything they know relevant from the discussion that day. The SI Leader can look at their work to assess the students' understanding of the course material and have them compare it to each other. It is important to at least discuss with the students so the SI Leader knows the students' comfort levels with the material. 
\end{enumerate}
\\ \hline
\end{tabular}

\newpage
\noindent
{\large\textbf{Instructions (please provide a \hl{DETAILED} activity breakdown below (and/or attached}}

\begin{tabular}{|p{0.9\linewidth}|}
\hline
\begin{enumerate}[label=\alph*.][itemsep=0.3em]
   \item Students are put into small groups and are asked to write down their confidence on a single notecard without names.
   \item The cards are evidence of the group's confidence in the content.
\end{enumerate}
\\ \hline
\end{tabular}

\vspace{0.5cm}
\noindent
{\large\textbf{Content (fully worked out problems \& answers)}

\begin{tabular}{|p{0.9\linewidth}|}
\hline
\begin{itemize}
   \item No new worked problems; brain dump.
\end{itemize}
\\ \hline
\end{tabular}

\vspace{0.5cm}
\noindent
{\large\textbf{Checking for Understanding (questions \& answers) }

\begin{tabular}{|p{0.9\linewidth}|}
\hline
\begin{itemize}
   \item Review confidence levels
   \item Interpret the variety of confidence and their distribution in groups.
   \item Get feedback for what they like and don't like
   \item Check in if the time is suitable for the students
\end{itemize}
\\ \hline
\end{tabular}
\newpage
\section*{Plan Checklist}

\begin{tabular}{|p{0.9\linewidth}|}
\hline
\begin{itemize}
\item[$\Box$] Variety of Collaborative Learning Techniques and Learning Strategies
\item[$\Box$] Chapter/slide\# where information is coming from
\item[$\Box$] Included timestamps (and buffer time as needed)
\item[$\Box$] Details on how leader will break up students
\item[$\Box$] Checking for Understanding questions
\item[$\Box$] Fully worked out problems and examples for each activity
\item[$\Box$] Necessary links required for online implementation (if applicable)
\item[$\Box$] Source material correctly cited (if applicable)
\item[$\Box$] Plan is uploaded to Teams
\end{itemize}
\\ \hline
\end{tabular}
\newpage

\end{document}